\paragraph{stable、奇维有理同伦、boundary 二挠与 exact-depth Haar 的四层分裂}
把定理 \ref{thm:conclusion-foldbin-stable-collapse-ordered-k0-memory}、定理 \ref{thm:conclusion-window6-rational-homotopy-inverts-tail-trichotomy}、推论 \ref{cor:conclusion-window6-rational-homotopy-determines-gauge-volume}、定理 \ref{thm:conclusion-window6-boundary-parity-directsummand-rational-blindness}、定理 \ref{thm:conclusion-window6-boundary-parity-three-layer-blind-filtration}、定理 \ref{thm:conclusion-window6-boundary-z6-torsor-local-global-mismatch}、定理 \ref{thm:conclusion-window6-groupoid-rather-than-free-quotient}、推论 \ref{cor:conclusion-window6-operational-commutative-faithful-triple-separation}、定理 \ref{thm:conclusion-binfold-gauge-complexity-three-scale-separation} 与推论 \ref{cor:conclusion-fibadic-uniform-divisor-randomization-gap} 并读，可把 window-\(6\) 与其 Fibonacci 深度平均进一步压缩为一条更紧的对象分类链。

\begin{theorem}[stable Morita 层与 order-unit 层的严格正交]\label{thm:conclusion-window6-stable-orderunit-strict-orthogonality}
设
\[
A_m:=\CC[\mathcal R_m^{\mathrm{bin}}]\cong \bigoplus_{w\in X_m}M_{d_w}(\CC),
\qquad
d_w:=d_m^{\mathrm{bin}}(w).
\]
则稳定层与单位层分别满足
\[
A_m\otimes\mathcal K\cong \mathcal K^{\oplus F_{m+2}},
\]
以及
\[
\bigl(K_0(A_m),K_0(A_m)^+,[1_{A_m}]\bigr)
\cong
\bigl(\ZZ^{X_m},\NN^{X_m},(d_w)_{w\in X_m}\bigr).
\]
因此，stable Morita 型只记住稳定类型总数 \(F_{m+2}\)，而全部块尺度谱 \((d_w)_{w\in X_m}\) 只能经由带单位元类的有序 \(K_0\) 被恢复；任何仅因子化于稳定同构类的不变量，都对纤维多重度分布严格失明。
\end{theorem}
\begin{proof}
这正是定理 \ref{thm:conclusion-foldbin-stable-collapse-ordered-k0-memory} 的直接重述。前一式给出稳定 Morita 层的完全塌缩；后一式表明 order unit \([1_{A_m}]\) 的坐标逐点记录了每个 Wedderburn 块的矩阵尺度。故两层所记忆的信息严格正交。证毕。
\end{proof}

\begin{theorem}[window-\texorpdfstring{$6$}{6} 连通有理包络对离散块谱与规范体积的完全反演]\label{thm:conclusion-window6-rational-envelope-reconstructs-block-spectrum}
\leanverified{paper\_conclusion\_window6\_rational\_envelope\_reconstructs\_block\_spectrum}
记
\[
A_6:=\CC[\mathcal R_6^{\mathrm{bin}}],
\qquad
K_6:=\Aut_\ast(A_6)^0
\cong
PU(2)^8\times PU(3)^4\times PU(4)^9,
\]
并设
\[
r_j:=\operatorname{rank}_{\QQ}\bigl(\pi_j(K_6)\otimes\QQ\bigr).
\]
则
\[
r_3=21,\qquad r_5=13,\qquad r_7=9,
\]
并且纤维直方图由三角反演唯一恢复：
\[
n_2=r_3-r_5=8,
\qquad
n_3=r_5-r_7=4,
\qquad
n_4=r_7=9.
\]
因此
\[
A_6\cong M_2(\CC)^{\oplus 8}\oplus M_3(\CC)^{\oplus 4}\oplus M_4(\CC)^{\oplus 9},
\]
\[
G_6^{\mathrm{bin}}\cong \mathfrak S_2^{\,8}\times \mathfrak S_3^{\,4}\times \mathfrak S_4^{\,9},
\]
并且
\[
\log |G_6^{\mathrm{bin}}|=39\log 2+13\log 3,
\qquad
g_6=\frac{39}{64}\log 2+\frac{13}{64}\log 3,
\qquad
\kappa_6=\frac{11}{8}\log 2+\frac{3}{16}\log 3.
\]
换言之，window-\(6\) 的连通有理同伦包络已经足以反演全部离散块谱与规范体积常数。
\end{theorem}
\begin{proof}
由定理 \ref{thm:conclusion-window6-rational-homotopy-inverts-tail-trichotomy}，
\[
(r_3,r_5,r_7)=(21,13,9)
\]
唯一恢复
\[
(n_2,n_3,n_4)=(8,4,9).
\]
再由推论 \ref{cor:conclusion-window6-rational-homotopy-determines-gauge-volume}，Wedderburn 分解、规范群直积分解以及 \(\log|G_6^{\mathrm{bin}}|\)、\(g_6\)、\(\kappa_6\) 的显式公式都由该直方图唯一确定。证毕。
\end{proof}

\begin{theorem}[有理连续完备性与 boundary 二挠盲性的并存]\label{thm:conclusion-window6-rational-completeness-with-boundary-torsion-blindness}
\leanverified{paper\_conclusion\_window6\_rational\_completeness\_with\_boundary\_torsion\_blindness}
在定理 \ref{thm:conclusion-window6-rational-envelope-reconstructs-block-spectrum} 的反演完备性成立的同时，window-\(6\) 的 boundary parity 仍形成一个与连通有理包络严格正交的离散层。更精确地，存在规范直和分解
\[
\bigl(G_6^{\mathrm{bin}}\bigr)^{\mathrm{ab}}
\cong
C_{\partial}\oplus C_{\mathrm{bulk}},
\qquad
C_{\partial}\cong (\ZZ/2\ZZ)^3,
\]
并且存在严格滤过
\[
0\subset P_6^{\mathrm{geo}}\subset P_6=C_{\partial},
\qquad
P_6^{\mathrm{geo}}\cong \ZZ/2\ZZ,
\qquad
P_6/P_6^{\mathrm{geo}}\cong (\ZZ/2\ZZ)^2.
\]
对任意连通 Lie 群 \(H\) 与任意同态
\[
\rho:\bigl(G_6^{\mathrm{bin}}\bigr)^{\mathrm{ab}}\to H,
\]
均有
\[
(\rho|_{C_{\partial}})^\ast:
H^{>0}(BH;\QQ)\longrightarrow H^{>0}(BC_{\partial};\QQ)=0.
\]
因此：
\begin{enumerate}
\normalfont
  \item 连通有理特征类对全部三位 boundary parity 严格全盲；
  \item 强局域几何 uplift 至多看见对角子群
  \[
  P_6^{\mathrm{geo}}\cong \ZZ/2\ZZ;
  \]
  \item 剩余两位 \(P_6/P_6^{\mathrm{geo}}\) 只能以纯离散 \(2\)-挠方式存在。
\end{enumerate}
\end{theorem}
\begin{proof}
规范直和分解与有理盲性来自定理 \ref{thm:conclusion-window6-boundary-parity-directsummand-rational-blindness}。严格滤过与几何对角可见子群来自定理 \ref{thm:conclusion-window6-boundary-parity-three-layer-blind-filtration}。把两者合并即可得到所述三项结论：连通有理包络可完全恢复块谱，却对 boundary 的三位 \(2\)-挠方向全部失明；强局域几何只保留同步对角 bit；其余两位则既非有理也非几何。证毕。
\end{proof}

\begin{theorem}[局域 torsor、全局群胚与三资源阈值的严格分裂]\label{thm:conclusion-window6-local-torsor-global-groupoid-three-resource-split}
\leanverified{paper\_conclusion\_window6\_local\_torsor\_global\_groupoid\_three\_resource\_split}
window-\(6\) 的边界升格扇区存在一个正则 \(\ZZ_6\)-torsor，
但其全窗口纤维划分既不可能由任何有限群自由作用的轨道实现，也不可能由任何阿贝尔群上的仿射陪集划分实现。于是，window-\(6\) 的精确对称承载对象必为群胚型半单对象，而不能属于等势自由商范畴；其自然的最小精确半单承载体恰为
\[
A_6\cong M_2(\CC)^{\oplus 8}\oplus M_3(\CC)^{\oplus 4}\oplus M_4(\CC)^{\oplus 9}.
\]
与此同时，操作可逆、交换可见与非交换忠实三条资源轴满足严格分离
\[
B_6^\star=2,
\qquad
\mu_{\mathrm{comm}}(A_6)=21,
\qquad
\mu_{\mathrm{faith}}(A_6)=64,
\]
从而
\[
2<21<64.
\]
因此，局域自由对称、交换分离预算与忠实非交换承载不是同一复杂度量的不同近似，而是三条互不代偿的资源轴。
\end{theorem}
\begin{proof}
边界升格扇区上的正则 \(\ZZ_6\)-torsor 与其不能向全窗口自由延拓，正是定理 \ref{thm:conclusion-window6-boundary-z6-torsor-local-global-mismatch} 的内容；而全窗口不能由等势自由商承载、只能落在群胚型半单对象上，则由定理 \ref{thm:conclusion-window6-groupoid-rather-than-free-quotient} 给出。

另一方面，推论 \ref{cor:conclusion-window6-operational-commutative-faithful-triple-separation} 已给出
\[
B_6^\star=2,\qquad \mu_{\mathrm{comm}}(A_6)=21,\qquad \mu_{\mathrm{faith}}(A_6)=64.
\]
故局域自由 torsor、交换可见、非交换忠实三条轴严格分离。证毕。
\end{proof}

\begin{theorem}[表示复杂度三尺度分离与 bounded exact-depth 平均的统一收缩律]\label{thm:conclusion-window6-representation-scales-vs-bounded-depth-contraction}
记
\[
\Xi_m:=\log k(G_m),
\qquad
\Gamma_m:=-\log \mathrm{cp}(G_m),
\qquad
\Sigma_m:=\log_2|H_2(G_m,\ZZ)|.
\]
则存在显式常数 \(C_{\mathrm{conj}},C_{\mathrm{vol}}>0\)，使得
\[
\Xi_m
=
C_{\mathrm{conj}}(2\varphi)^{m/2}
+o\!\bigl((2\varphi)^{m/2}\bigr),
\]
\[
\Gamma_m
=
2^m\left(m\log\frac{2}{\varphi}-C_{\mathrm{vol}}\right)
-C_{\mathrm{conj}}(2\varphi)^{m/2}
+o\!\bigl((2\varphi)^{m/2}\bigr),
\]
并且当 \(m\) 充分大时
\[
H_2(G_m,\ZZ)\cong (\ZZ/2\ZZ)^{\,n_m(n_m+1)/2},
\qquad
\Sigma_m=\frac12\,n_m(n_m+1),
\qquad
n_m=F_{m+2}.
\]
因此三者分别驻留在
\[
(2\varphi)^{m/2},\qquad 2^m,\qquad \varphi^{2m}
\]
三条严格分离的尺度上，不存在统一归一化。

另一方面，对任意
\[
M=\prod_{j=1}^{r}p_j^{a_j},
\qquad
A_M:=\frac1{\tau(M)}\sum_{d\mid M}E_d,
\]
在每个 exact-depth 新空间 \(\mathcal N_n\) 上都有
\[
A_M|_{\mathcal N_n}
=
\begin{cases}
\dfrac{\tau(M/n)}{\tau(M)}\,\mathrm{id}, & n\mid M,\\[1ex]
0, & n\nmid M.
\end{cases}
\]
并且若一族整数 \(M_\nu\) 满足
\[
\sup_{\nu,j}a_j^{(\nu)}\le A<\infty,
\]
则存在统一常数
\[
\varepsilon=\frac1{A+1}>0
\]
使得
\[
\|A_{M_\nu}^k-E_1\|_{2\to 2}\le (1-\varepsilon)^k
\qquad(\forall \nu,\ k\ge 1).
\]
故一切 bounded-exponent 的除数随机化都不可能在非平凡 exact-depth 扇区上逼近恒等；它们只能以统一指数速率收缩到常数投影。
\end{theorem}
\begin{proof}
前三条渐近与三尺度分离正是定理 \ref{thm:conclusion-binfold-gauge-complexity-three-scale-separation} 的内容。对 random-depth 平均，推论 \ref{cor:conclusion-fibadic-uniform-divisor-randomization-gap} 已给出
\[
A_M|_{\mathcal N_n}
=
\begin{cases}
\dfrac{\tau(M/n)}{\tau(M)}\,\mathrm{id}, & n\mid M,\\[1ex]
0, & n\nmid M,
\end{cases}
\qquad
1-\rho_2(A_M)=\min_j\frac1{a_j+1}.
\]
若 \(\sup_{\nu,j}a_j^{(\nu)}\le A\)，则
\[
\rho_2(A_{M_\nu})\le \frac{A}{A+1}=1-\varepsilon.
\]
再由定理 \ref{thm:conclusion-fibadic-random-depth-average-spectral-criterion} 的精确范数公式
\[
\|A_M^k-E_1\|_{2\to 2}
=
\sup_{n>1:\,\mathcal N_n\neq 0}\widehat\alpha(n)^k,
\]
立即得到所示统一收缩率。证毕。
\end{proof}

\begin{conjecture}[window-\texorpdfstring{$6$}{6} 忠实承载的三寄存器与四寄存器原理]\label{conj:conclusion-window6-three-four-register-faithful-carrier}
若一个对象要对 window-\(6\) 的完整层级保持 faithful，则至少应同时外置如下三类彼此正交的寄存器：
\[
\mathfrak R_{\mathrm{st}}
\quad\text{（稳定块数层，承载 }F_{m+2}\text{）},
\]
\[
\mathfrak R_{\mathrm{ord}}
\quad\text{（order-unit / odd-sphere 层，承载 }(d_w)\text{ 或等价地 }(r_3,r_5,r_7)\text{）},
\]
\[
\mathfrak R_{\partial}
\quad\text{（boundary 二挠层，承载 }0\subset P_6^{\mathrm{geo}}\subset P_6\text{）}.
\]
若进一步要求与 Fibonacci 深度平均的可审计性兼容，则还必须附加一条 commuting depth register，
即由
\[
\{\nu_d\}_{d\ge 1}
\]
生成的 exact-depth signed Haar 体系。换言之，faithful 对象至少应具有
\[
\mathfrak R_{\mathrm{st}}\oplus \mathfrak R_{\mathrm{ord}}\oplus \mathfrak R_{\partial}
\]
的三层结构；若再要求跨深度审计，则应扩张为四层。
\end{conjecture}
